\section{\personFloer{}边界映射的定义}
\label{sec:def-hamiltonian-floer-boundary}

设 $H$ 是非退化的一周期哈密顿量，$J$ 是相容殆复结构的一周期族。我们先定义\personFloer{}边界映射，并指出定义这一对的\personFloer{}同调 $HF_*(H,J)$ 所需的横截性条件。
在 $(M,\omega)$ 的半正性假设下，我们将完整讨论这个横截性问题。

\begin{definition}
\label{def:relative-cylinder-homotopy-class}
设 $z,z'\in\operatorname{Per}(H)$。以 $\pi_2(z,z')$ 表示相对于边界
\[
 u(0,t)=z(t),\qquad u(1,t)=z'(t)
\]
的光滑映射
\[
 u:[0,1]\times S^1=:T\longrightarrow M
\]
的同伦类集合，并以 $[u]\in\pi_2(z,z')$ 表示 $u$ 的同伦类。
\end{definition}

类似地，以 $\pi_2(z)$ 表示映射
\[
 w:D^2\longrightarrow M;\qquad w|_{\partial D^2}=z
\]
的相对同伦类集合。注意，$\pi_2(M)$ 在 $\pi_2(z)$ 和 $\pi_2(z,z')$ 上有自然的
“粘合一个球面”作用（模去 $\pi_1(M)$ 的作用）。此外，对
$C\in\pi_2(z,z')$，存在由粘合映射 $w\mapsto w\#u$ 诱导的自然映射
\[
 (\,\cdot\,)\#C:\pi_2(z)\longrightarrow\pi_2(z').
\]
具体地说，我们固定在 $D^2$ 的极坐标 $(r,\theta)$ 中按下式定义
$w\#u:D^2\to M$：
\begin{equation}
 (w\#u)(r,\theta)=
 \begin{cases}
  w(2r,\theta),&0\leq r\leq\frac12,\\
  u(2r-1,\theta),&\frac12\leq r\leq1.
 \end{cases}
\label{eq:19-disc-cylinder-gluing}
\end{equation}
还有自然的粘合映射
\[
 \pi_2(z_0,z_1)\times\pi_2(z_1,z_2)\longrightarrow\pi_2(z_0,z_2),
 \qquad (u_1,u_2)\longmapsto u_1\#u_2.
\]

\stfhIndex{I167}{XiangduiKangli-Zengdeerzhibiao}{相对康利--曾德尔指标（relative Conley--Zehnder index）}
现在定义\personConley{}--\personZehnder{}指标
$\mu_H:\pi_2(z,z')\to\mathbb Z$。只要 $z_0,z_1$ 位于 $\Omega(M)$ 的同一连通分支，
这个定义并不要求它们可缩。

令 $T=[0,1]\times S^1$。对给定映射 $u:T\to M$，选取一个辛平凡化
\[
 \Phi:u^*TM\longrightarrow T\times\mathbb R^{2n}.
\]
要求它满足标记条件
\begin{equation}
 \Phi(t,1)\equiv\Phi(0,1),\qquad t\in[0,1].
\label{eq:19-trivialization-marking}
\end{equation}
我们知道，对某些 $p_0,p_1\in\operatorname{Fix}\phi_H^1$，有
$z_0(t)=\phi_H^t(p_0)$、$z_1(t)=\phi_H^t(p_1)$。于是有两个映射
\[
 \alpha_{\Phi,i}:[0,1]\longrightarrow\operatorname{Sp}(2n),\qquad i=0,1,
\]
使得对 $v\in\mathbb R^{2n}$ 和 $t\in[0,1]$，
\[
 \Phi\circ d\phi_H^t(p_i)\circ\Phi^{-1}(i,t,v)
 =(i,t,\alpha_{\Phi,i}(t)v).
\]
由条件~\eqref{eq:19-trivialization-marking} 及 $H$ 的非退化性，对 $i=0,1$ 都有
\[
 \alpha_{\Phi,i}(0)=\operatorname{id},\qquad
 \alpha_{\Phi,i}(1)\in\operatorname{Sp}^*(2n),
\]
因而\personConley{}--\personZehnder{}指标 $\mu_{CZ}(\alpha_{\Phi,i})$ 有定义。

\begin{definition}
\label{def:relative-conley-zehnder-index}
对每个 $C\in\pi_2(z,z')$，定义其相对\personConley{}--\personZehnder{}指标为
\begin{equation}
 \mu_H(z,z';C):=\mu_{CZ}(\alpha_{\Phi,0})-
 \mu_{CZ}(\alpha_{\Phi,1}).
\label{eq:19-relative-cz-index}
\end{equation}
边界条件不会引起混淆时，也写作 $\mu_H(C)$。
\end{definition}

\begin{exercise}
\label{ex:relative-cz-trivialization-independent}
证明此定义与满足~\eqref{eq:19-trivialization-marking} 的平凡化 $\Phi$ 无关。
\end{exercise}

下面的命题把这个相对指标与
$[z,w]\in\operatorname{Crit}\mathcal A_H$ 的绝对\personConley{}--\personZehnder{}指标
$\mu_H([z,w])$ 联系起来。

\begin{proposition}
\label{prop:relative-absolute-cz-relation}
对类 $C\in\pi_2(z,z')$ 的\emph{一个（因而任意一个）代表}
$u:[0,1]\times S^1\to M$
\footnote{译者注：原书在代表映射的类型中印作
$u:[0,1]\times S^1\times M$，缺少映射箭头；结合本节定义及随后证明，译文补为
$u:[0,1]\times S^1\to M$。}
，有
\begin{equation}
 \mu_H(z,z';C)=\mu_H([z,w])-\mu_H([z',w\#u]).
\label{eq:19-relative-absolute-cz}
\end{equation}
\end{proposition}

\begin{proof}
选取平凡化
\[
 \Psi_w:w^*TM\longrightarrow D^2\times\mathbb R^{2n}
\]
以及满足~\eqref{eq:19-trivialization-marking} 和
\begin{equation}
 \Psi_w|_{\partial D^2}=\Phi|_{\{0\}\times S^1}
\label{eq:19-compatible-glued-trivializations}
\end{equation}
的平凡化 $\Phi:u^*TM\to T\times\mathbb R^{2n}$；这里把
$\partial D^2\cong\{0\}\times S^1$。粘合 $\Psi_w\#\Phi$ 自然诱导
$(w\#u)^*TM$ 的平凡化。由构造和条件~\eqref{eq:19-compatible-glued-trivializations}，
\[
 \mu_H(\alpha_{\Psi_w\#\Phi})=\mu_H(\alpha_{\Phi,1}),
 \qquad
 \mu_H(\alpha_{\Psi_w})=\mu_H(\alpha_{\Phi,0}).
\]
因此
\[
\begin{aligned}
 \mu_H([z',w\#u])
 &=\mu_H(\alpha_{\Psi_w\#\Phi})\\
 &=\mu_H(\alpha_{\Phi,1})-\mu_H(\alpha_{\Phi,0})
   +\mu_H(\alpha_{\Phi,0})\\
 &=-\mu_H(z,z';[u])+\mu_H([z,w]),
\end{aligned}
\]
证明完毕。
\footnote{译者注：原书本计算首项漏印右方括号，写成
$\mu_H([z',w\#u]$；译文补齐。}
\end{proof}

现在以
\[
 \widetilde{\mathcal M}(H,J;z,z';C)
\]
表示满足渐近条件与同伦条件
\begin{equation}
 u(-\infty)=z,\qquad u(\infty)=z';\qquad [u]=C
\label{eq:19-floer-asymptotic-homotopy}
\end{equation}
的方程 (18.5.41) 的有限能量解集，并令
\[
 \mathcal M(H,J;z,z';C)
 =\widetilde{\mathcal M}(H,J;z,z';C)/\mathbb R.
\]
这里要指出，虽然 $u$ 定义在 $\mathbb R\times S^1$ 上，但由于方程 (18.5.41) 的
解具有指数衰减性质（回忆我们假设 $H$ 非退化），它可紧化为满足相应边界条件
\[
 \bar u(0)=z,\qquad\bar u(1)=z'
\]
的连续映射 $\bar u:[0,1]\times S^1\to M$。称 $\bar u$ 为 $u$ 的\emph{紧化映射}。
略用记号，也以 $[u]$ 表示紧化映射的类 $[\bar u]\in\pi_2(z,z')$。

下面是\personFredholm{}指标的基本公式。（证明见 (Schw95)。）

\begin{proposition}
\label{prop:closed-floer-index-formula}
设 $u\in\mathcal M(H,J;z,z';C)$，则
\begin{equation}
 \operatorname{Index}u=\mu_H(z,z';C).
\label{eq:19-closed-floer-index}
\end{equation}
特别地，若 $J$ 是 $H$-正则的，则
\[
 \dim\widetilde{\mathcal M}(H,J;z,z';C)=\mu_H(z,z';C).
\]
\end{proposition}

\begin{exercise}
\label{ex:closed-string-index-formula}
把命题 15.3.1 的证明转写到闭弦情形，并证明此命题。
\end{exercise}

\personFloer{}边界映射
\[
 \partial_{(H,J)}:CF_{k+1}(H)\longrightarrow CF_k(H)
\]
通过计数满足 $\mu_H(z,z';C)=1$ 的三元组 $(z,z';C)$ 所对应的
$\mathcal M(H,J;z,z';C)$ 中的元素来定义，条件如下。

\begin{definition}[边界映射]
\label{def:hamiltonian-floer-boundary-regularity}
设 $H$ 非退化。假设 $J$ 满足：
\begin{enumerate}[label=(\arabic*)]
\item 对任意满足
\[
 \mu_H(z_0,z_1;C)
 =\mu_H([z_0,w_0])-\mu_H([z_1,w_0\#C])\leq0
\]
的对子 $z_0,z_1\in\operatorname{Per}(H)$，除非 $z_0=z_1$ 且 $C=0$，否则
$\widetilde{\mathcal M}(H,J;z_0,z_1;C)=\varnothing$。当 $z_0=z_1$ 且 $C=0$ 时，
仅有定常解，即对所有 $\tau\in\mathbb R$ 都有 $u(\tau)\equiv z_0=z_1$。
\item 对任意 $z_0,z_1\in\operatorname{Per}(H)$ 以及满足
\[
 \mu_H(z_0,z_1;C)=1
\]
的同伦类 $C\in\pi_2(z_0,z_1)$，空间
$\widetilde{\mathcal M}(H,J;z_0,z_1;C)$ 横截，且商
\[
 \mathcal M(H,J;z_0,z_1;C)
 =\widetilde{\mathcal M}(H,J;z_0,z_1;C)/\mathbb R
\]
紧致，因而是有限集。记
\[
 n(H,J;z_0,z_1;C)=\#\mathcal M(H,J;z_0,z_1;C)
\]
为此空间中元素的代数计数。若三元组 $(z,z';C)$ 不满足指标要求，则约定
\[
 n(H,J;z_0,z_1;C)=0.
\]
\item 对任意 $z_0,z_2\in\operatorname{Per}(H)$ 和满足
\[
 \mu_H(z_0,z_2;C)=2
\]
的 $C\in\pi_2(z_0,z_2)$，$\mathcal M(H,J;z_0,z_2;C)$ 可紧化为光滑一维流形，
其边界由所有可能的断裂流线
\[
 [u_1]\#_\infty[u_2]
\]
组成，其中对某个 $y\in\operatorname{Per}(H)$ 以及
$C_1\in\pi_2(z_0,y)$、$C_2\in\pi_2(y,z_2)$，
\[
 u_1\in\widetilde{\mathcal M}(H,J;z_0,y;C_1),\qquad
 u_2\in\widetilde{\mathcal M}(H,J;y,z_2;C_2),
\]
并且
\[
 C_1\#C_2=C,\qquad [u_1]=C_1,\quad[u_2]=C_2,
\]
以及
\[
 \mu_H(z_0,y;C_1)=\mu_H(y,z_2;C_2)=1.
\]
这里以 $[u]$ 表示由 $u$ 代表的等价类。
\end{enumerate}
\stfhIndex{I089}{Fuluoerzhengze}{弗洛尔正则（Floer regular）}
称任何这样的 $J$ 为 \emph{$H$-正则}，并称这样的对子 $(H,J)$ 为\emph{\personFloer{}正则}。
\end{definition}

这个定义的要点是：对\personFloer{}正则对子 $(H,J)$，\personFloer{}边界映射
\[
 \partial=\partial_{(H,J)}:CF_*(H)\longrightarrow CF_*(H)
\]
有定义且满足 $\partial\partial=0$，于是可以取其同调。由于 $CF_*(H)$ 由
$[z,w]\in\operatorname{Crit}\mathcal A_H$ 生成，我们需要更详细地说明这个构造。

给定 $[z,w]$。对每个 $[z',w']\in\operatorname{Crit}\mathcal A_H$，考虑所有满足
\begin{equation}
 [z',w\#C]=[z',w'],\qquad\mu_H(z,z';C)=1
\label{eq:19-boundary-class-condition}
\end{equation}
的 $C\in\pi_2(z,z')$，并定义模空间
\[
 \mathcal M(H,J;[z,w],[z',w'])
 :=\bigcup_C\{\mathcal M(H,J;z,z';C)\mid
 C\text{ 满足~\eqref{eq:19-boundary-class-condition}}\}.
\]
按 $\Gamma$-覆盖空间 $\widetilde{\mathcal L}_0(M)$ 的定义，可能有不止一个
$C\in\pi_2(z^-,z^+)$ 满足~\eqref{eq:19-boundary-class-condition}。不过，\personGromov{}--\personFloer{}紧性给出如下有限性。

\begin{lemma}
\label{lem:finitely-many-boundary-classes}
上述并是有限并。换言之，对任意给定对子 $([z,w],[z',w'])$，只有有限多个
$C\in\pi_2(z,z')$ 同时满足~\eqref{eq:19-boundary-class-condition} 和
$\mathcal M(H,J;z,z';C)\neq\varnothing$。
\end{lemma}

\begin{proof}
反设存在无穷序列
$u_i\in\mathcal M(H,J;[z,w],[z',w'])$，且 $[u_i]\in\pi_2(z,z')$ 两两不同。
首先注意
\[
 E_{(H,J)}(u_i)=\mathcal A_H([z,w])-\mathcal A_H([z',w'])
\]
与 $i$ 无关。由\personGromov{}--\personFloer{}紧性，可取仍记为 $u_i$ 的子列，使其收敛到
\[
 u_\infty=u^1\#u^2\#\cdots\#u^N,
\]
其中每个 $u^k$ 形如
\[
 u^k=u_{k,0}+\sum_{\ell=1}^{L_k}v_\ell^k.
\]
这里，$\{u_{k,0}\}_k$ 构成一条从 $z$ 到 $z'$ 的连通\personFloer{}流线链，
$v_\ell^k$ 是附着于 $u_{k,0}$ 的球面气泡；每个 $u_{k,0}$ 是连接周期轨道
$z_{k-1},z_k$ 的光滑\personFloer{}流线，其中 $z_0=z$、$z_N=z'$。

此外，\personGromov{}--\personFloer{}紧性还蕴含：对所有充分大的 $i$，有与 $i$ 无关的等式
\footnote{译者注：原书此处只写
$[u_i]=[u_\infty]+\sum_{\ell=1}^{L_k}[v_\ell^k]$，其中 $k$ 未被量化；而上文
$u_\infty$ 含全部 $N$ 个分支。译文据该分解补出对 $k$ 的求和。}
\[
 [u_i]=[u_\infty]+\sum_{k=1}^{N}\sum_{\ell=1}^{L_k}[v_\ell^k].
\]
这与 $[u_i]\in\pi_2(z,z')$ 两两不同的假设矛盾。
\end{proof}

现在把 $u$ 看成覆盖空间 $\widetilde{\mathcal L}_0(M)$ 中的道路，把
$u\in\mathcal M(H,J;[z,w],[z',w'])$ 的渐近条件写成
\begin{equation}
 u(-\infty)=[z,w],\qquad u(\infty)=[z',w'].
\label{eq:19-lifted-asymptotic-condition}
\end{equation}
\personFloer{}边界映射
$\partial=\partial_{(H,J)}:CF_*(H)\to CF_*(H)$ 由矩阵系数
\[
 \langle\partial([z,w]),[z',w']\rangle
 :=\sum_C n_{(H,J)}(z,z';C)
\]
定义，其中 $C$ 如~\eqref{eq:19-boundary-class-condition}，且 $[z,w]$ 与 $[z',w']$
的\personConley{}--\personZehnder{}指标满足
\[
 \mu_H([z,w])-\mu_H([z',w'])=\mu(z,z';C)=1.
\]
由定义 19.2.7 的定义与假设，$\partial=\partial_{(H,J)}$ 的次数为 $-1$，且
$\partial\circ\partial=0$。

\stfhIndex{I084}{Fuluoerbilian}{弗洛尔闭链（Floer cycle）}
\begin{definition}
\label{def:floer-cycle-boundary-homology}
若\personFloer{}链 $\beta\in CF(H)$ 满足 $\partial\beta=0$，即
$\beta\in\ker\partial_{(H,J)}$，则称 $\beta$ 是 $(H,J)$ 的\emph{\personFloer{}闭链}；若
$\beta\in\operatorname{Image}\partial_{(H,J)}$，则称它为\emph{\personFloer{}边界}。若
$\beta'-\beta$ 是边界，则称两个\personFloer{}链 $\beta,\beta'$ \emph{同调}。
\end{definition}

定义 $(H,J)$ 的\personFloer{}同调为
\[
 HF_*(H,J):=\ker\partial_{(H,J)}/\operatorname{Image}\partial_{(H,J)}.
\]
它既可视为分次 $\mathbb Q$-向量空间，也可视为 $\Lambda_\omega$-模。
